\documentclass[11pt]{amsart}

\usepackage[T1]{fontenc}
\usepackage[utf8]{inputenc}
\usepackage{lmodern}
\usepackage[margin=1.08in]{geometry}
\usepackage{microtype}
\usepackage{amsmath,amssymb,amsthm,mathtools,mathrsfs}
\usepackage{enumitem}
\usepackage{xcolor}
\usepackage[colorlinks=true,linkcolor=blue!50!black,citecolor=blue!50!black,urlcolor=blue!50!black]{hyperref}

\numberwithin{equation}{section}

\newtheorem{theorem}{Theorem}[section]
\newtheorem{proposition}[theorem]{Proposition}
\newtheorem{corollary}[theorem]{Corollary}
\newtheorem{lemma}[theorem]{Lemma}
\theoremstyle{definition}
\newtheorem{definition}[theorem]{Definition}
\theoremstyle{remark}
\newtheorem{remark}[theorem]{Remark}

\DeclareMathOperator{\Arg}{Arg}
\DeclareMathOperator{\Fix}{Fix}
\DeclareMathOperator{\ord}{ord}
\DeclareMathOperator{\Rea}{Re}
\DeclareMathOperator{\Ima}{Im}

\newcommand{\ZZ}{\mathbb Z}
\newcommand{\RR}{\mathbb R}
\newcommand{\CC}{\mathbb C}
\newcommand{\NN}{\mathbb N}
\newcommand{\BB}{\mathbf B}
\newcommand{\PP}{\mathbb P}
\newcommand{\Sph}{\mathbb S^1}
\newcommand{\ph}{\operatorname{ph}}
\newcommand{\taup}{\tau}
\newcommand{\e}{\mathrm e}
\newcommand{\ii}{\mathrm i}
\newcommand{\cO}{\mathcal O}

\title[Atomic phase packets of zeta zeros]{Taylor Atomization, Orthogonal Walks, and Spectral Packets\ for the Non-Trivial Zeta Zeros}
\author{}
\date{May 2026}

\begin{document}
\begin{abstract}
For every zero $\rho$ of the Riemann zeta function one has $F_1(\rho)=\zeta(\rho)-1=-1$ in the globalization semiring $G(\ZZ)=\ZZ\sqcup\{\taup\}$, so the reduced split-zero coefficient is the fixed Boolean pair $A=\{0,\taup\}\cong\BB$.  The pointwise coefficient is constant, and the problem is therefore geometric.

In this paper we make precise the orthogonal-walk and atomic-phase formalism extracted from unpublished correspondence of Jacolm Tobley.  First, for a finite real sequence $(\alpha_1,\dots,\alpha_m)$ we prove the exact orthogonal-walk expansion
\[
\prod_{k=1}^m (1+\ii\alpha_k)=\sum_{r=0}^m \ii^r e_r(\alpha_1,\dots,\alpha_m),
\]
with step lengths given by the elementary symmetric polynomials, and we identify the normalized product with $\exp\bigl(\ii\sum_k \arctan\alpha_k\bigr)$.  Second, given a real-analytic phase series $\Theta(t)=\sum_{n\ge 1}\beta_n t^n$, we prove that on a sufficiently small interval the exact factorization
\[
\e^{\ii\Theta(t)}=\prod_{n\ge 1}\frac{1+\ii\tan(\beta_n t^n)}{\sqrt{1+\tan^2(\beta_n t^n)}}
\]
holds, and that the corresponding partial products are normalized endpoints of finite orthogonal walks.  Thus Taylor phase series become precise complex exponentials built from atomic sequences.

Third, for a non-trivial zero $\rho$ of the completed zeta function $\xi$ of multiplicity $m$, we apply this formalism to the reduced local germ
\[
\psi_\rho(z):=\frac{\xi(\rho+z)}{\lambda_\rho z^m},\qquad \lambda_\rho:=\frac{\xi^{(m)}(\rho)}{m!},
\]
and obtain, for every unit direction $\nu\in\Sph$, a canonical real-analytic local phase packet
\[
\ph\bigl(\xi(\rho+t\nu)\bigr)=\nu^m\ph(\lambda_\rho)\exp\Bigl(\ii\sum_{n\ge 1}\beta_{\rho,\nu,n}t^n\Bigr)
\qquad (0<t\ll 1).
\]
The finite truncations are normalized endpoints of explicit orthogonal walks determined by the atomic slopes $\tan(\beta_{\rho,\nu,n}t^n)$.  Finally, we combine this local phase formalism with the affine spectral coordinate $\Phi(s)=\ii(s-\tfrac12)$ and the Connes--Consani $\zeta$-cycle lengths.  In this way a non-trivial zeta zero is represented geometrically by a weighted Boolean spectral packet together with a canonical Taylor-atomic phase sequence.
\end{abstract}

\maketitle

\tableofcontents

\section{Introduction}

Let
\[
S=G(\ZZ)=\ZZ\sqcup\{\taup\},
\qquad
A:=\{0,\taup\}\subset S,
\qquad
F_1(s):=\zeta(s)-1.
\]
If $\rho$ is a zero of $\zeta$, then $F_1(\rho)=-1$, and multiplication by $-1$ fixes exactly the Boolean pair $A\cong\BB$.  Thus the reduced split-zero coefficient attached to a zeta zero is constant.  The mathematical issue is not pointwise evaluation but the extraction of geometric data from the local analytic germ and from the global packet symmetries.

The present paper addresses a specific question.  The previous split-zero papers produced divisor sheaves, packet sheaves, jet enhancements, and compactified tails.  What they did not yet do is to convert the local Taylor expansions at a zero into a precise complex-exponential formalism governed by the orthogonal-walk constructions from the unpublished correspondence of Jacolm Tobley \cite{Tobley}.  The point of the present note is to do exactly that.

There are three inputs.
\begin{enumerate}[label=(\roman*)]
\item The split-zero coefficient object is the Boolean pair $A=\{0,\taup\}\cong\BB$.
\item The non-trivial zero packets are governed by the completed zeta function
\[
\xi(s)=\frac12 s(s-1)\pi^{-s/2}\Gamma\!\left(\frac s2\right)\zeta(s),
\]
which satisfies $\xi(1-s)=\xi(s)$ and $\overline{\xi(\bar s)}=\xi(s)$.
\item Tobley's phase formalism rewrites products of normalized affine factors
\[
\frac{1+\ii\alpha}{\sqrt{1+\alpha^2}}
\]
as exact complex exponentials and expands the unnormalized products into orthogonal walks whose step lengths are elementary symmetric polynomials \cite{Tobley}.
\end{enumerate}

The new observation is that the local phase of the reduced zeta germ at a zero is exactly of the form treated by the Tobley formalism.  After dividing out the vanishing factor $z^m$ and the leading coefficient $\lambda_\rho$, one is left with a nonvanishing holomorphic germ.  A holomorphic logarithm of that germ has a Taylor expansion, and the imaginary part of that logarithm produces a real phase series
\[
\Theta_{\rho,\nu}(t)=\sum_{n\ge 1}\beta_{\rho,\nu,n}t^n.
\]
The corresponding local unit phase is $\exp(\ii\Theta_{\rho,\nu}(t))$.  The point is that this is not merely an abstract exponential: it can be written exactly as an infinite product of normalized affine factors and approximated by finite orthogonal walks with explicit step lengths.

This gives a concrete packet-geometric answer to the question of how a non-trivial zero maps into the split-zero formalism.  The reduced coefficient is still the Boolean pair $A$; the multiplicity contributes the Boolean cube $A^{m_\rho}$; the packet position is carried by the affine spectral coordinate
\[
\Phi(s)=\ii\Bigl(s-\frac12\Bigr);
\]
and the local phase is carried by a canonical Taylor-atomic sequence.  On the critical line, the same packet also determines the primitive $\zeta$-cycle length $2\pi/|\Im\rho|$ by the theorem of Connes and Consani \cite{ConnesConsaniCycles}.

\section{Split-zero coefficients and spectral packets}

\subsection{The Boolean coefficient object}

\begin{definition}
Let
\[
S:=\ZZ\sqcup\{\taup\},\qquad \taup\notin\ZZ,
\]
with operations
\[
m\oplus n:=m+n,
\qquad
m\oplus \taup=\taup\oplus m:=m,
\qquad
\taup\oplus\taup:=\taup,
\]
and
\[
m\odot n:=mn,
\qquad
m\odot\taup=\taup\odot m:=\taup,
\qquad
\taup\odot\taup:=\taup.
\]
Thus $\taup$ is the additive identity and multiplicative absorber, while $1\in\ZZ\subset S$ is the multiplicative identity.
\end{definition}

\begin{proposition}\label{prop:boolean-pair}
The subset
\[
A:=\{0,\taup\}\subset S
\]
is a subsemiring, and the map
\[
A\longrightarrow \BB=\{0,1\},\qquad \taup\mapsto 0,\qquad 0\mapsto 1,
\]
is an isomorphism of semirings.
\end{proposition}

\begin{proof}
Inside $A$ one has
\[
0\oplus0=0,
\qquad
0\oplus\taup=0,
\qquad
\taup\oplus\taup=\taup,
\]
and
\[
0\odot0=0,
\qquad
0\odot\taup=\taup,
\qquad
\taup\odot\taup=\taup.
\]
These are exactly the Boolean addition and multiplication tables under the displayed identification.
\end{proof}

\begin{proposition}\label{prop:F1-zero}
For every zero $\rho$ of $\zeta$ one has
\[
F_1(\rho)=-1.
\]
The unique order-$2$ unit of $S$ is $-1$, and its fixed-point set is exactly $A$.
\end{proposition}

\begin{proof}
If $\zeta(\rho)=0$, then $F_1(\rho)=\zeta(\rho)-1=-1$.  Since the only units of $\ZZ$ are $\pm1$, the only units of $S$ are also $\pm1$.  The element $-1$ has order $2$, while $1$ does not.  Finally,
\[
(-1)\odot\taup=\taup,
\qquad
(-1)\odot n=n\iff -n=n\iff n=0,
\]
so the fixed-point set is $\{0,\taup\}=A$.
\end{proof}

\subsection{Non-trivial packets and the spectral coordinate}

\begin{definition}
Let
\[
\xi(s):=\frac12 s(s-1)\pi^{-s/2}\Gamma\!\left(\frac s2\right)\zeta(s)
\]
be the completed zeta function.  Its zero set is denoted
\[
Z_\xi:=\{\rho\in\CC:\xi(\rho)=0\}.
\]
These are exactly the non-trivial zeros of $\zeta$, counted with the same multiplicities.
\end{definition}

\begin{proposition}\label{prop:packet-coordinate}
The set $Z_\xi$ is stable under complex conjugation and under $s\mapsto 1-s$.  If $\rho\in Z_\xi$ has multiplicity $m_\rho$, then every point in the orbit
\[
P_\rho:=\{\rho,\bar\rho,1-\rho,1-\bar\rho\}
\]
has the same multiplicity.  The affine map
\[
\Phi:\CC\longrightarrow \CC,
\qquad
\Phi(s):=\ii\Bigl(s-\frac12\Bigr),
\]
satisfies
\[
\Phi(1-s)=-\Phi(s),
\qquad
\Phi(\bar s)=-\overline{\Phi(s)}.
\]
Consequently,
\[
\Phi(P_\rho)=\{z,-z,\bar z,-\bar z\},
\qquad z:=\Phi(\rho).
\]
If $\rho=\tfrac12+\ii\gamma$ lies on the critical line, then $z=-\gamma\in\RR$.
\end{proposition}

\begin{proof}
The functional equation $\xi(1-s)=\xi(s)$ implies that $1-\rho$ is a zero whenever $\rho$ is.  The identity $\overline{\xi(\bar s)}=\xi(s)$ implies that $\bar\rho$ is a zero whenever $\rho$ is.  Because these maps are biholomorphic and preserve $\xi$ up to composition with a nowhere-vanishing factor of modulus $1$, the multiplicity is constant on each orbit.

The displayed formulas for $\Phi$ are immediate:
\[
\Phi(1-s)=\ii\Bigl(\frac12-s\Bigr)=-\ii\Bigl(s-\frac12\Bigr)=-\Phi(s),
\]
and
\[
\Phi(\bar s)=\ii\Bigl(\bar s-\frac12\Bigr)=-\overline{\ii\Bigl(s-\frac12\Bigr)}=-\overline{\Phi(s)}.
\]
Applying these to the orbit of $\rho$ gives the packet formula.  If $\rho=\tfrac12+\ii\gamma$, then
\[
\Phi(\rho)=\ii\bigl(\ii\gamma\bigr)=-\gamma\in\RR.
\]
\end{proof}

\begin{remark}
For a zero $\rho$ of multiplicity $m_\rho$, the multiplicity coefficient object in the split-zero formalism is the Boolean cube $A^{m_\rho}$.  At the reduced level one sees only the fixed Boolean pair $A$; the higher Boolean cube records multiplicity.
\end{remark}

\begin{corollary}[Critical packets and $\zeta$-cycle lengths]\label{cor:critical-length}
Let $\rho=\tfrac12+\ii\gamma\in Z_\xi$ with $\gamma\neq0$.  Then the packet coordinate is $\Phi(\rho)=-\gamma$, and the primitive $\zeta$-cycle length attached to the packet is
\[
\Lambda_\rho:=\frac{2\pi}{|\gamma|}=\frac{2\pi}{|\Phi(\rho)|}.
\]
\end{corollary}

\begin{proof}
The identity $\Phi(\rho)=-\gamma$ is Proposition~\ref{prop:packet-coordinate}.  Connes and Consani prove that if $\zeta(\tfrac12+\ii s)=0$ with $s>0$, then every circle of length an integral multiple of $2\pi/s$ is a $\zeta$-cycle, and that $s$ appears in the corresponding spectrum \cite[Theorem~1.1(ii)]{ConnesConsaniCycles}.  This yields the displayed primitive length.
\end{proof}

\section{Atomic phase factors and orthogonal walks}

The exact phase formalism used below is extracted from Tobley's unpublished correspondence on arctangent decompositions, orthogonal walks, and atomic step models of complex multiplication \cite{Tobley}.

\begin{definition}
For $a\in\RR$, define the normalized affine phase factor
\[
u(a):=\frac{1+\ii a}{\sqrt{1+a^2}}\in\Sph.
\]
For a finite sequence $a_1,\dots,a_M\in\RR$, define the unnormalized and normalized products
\[
U_M(a_1,\dots,a_M):=\prod_{n=1}^M (1+\ii a_n),
\qquad
\widehat U_M(a_1,\dots,a_M):=\prod_{n=1}^M \nu(a_n).
\]
\end{definition}

\begin{proposition}[Finite orthogonal-walk expansion]\label{prop:finite-orthogonal}
Let $a_1,\dots,a_M\in\RR$, and let $e_r(a_1,\dots,a_M)$ be the elementary symmetric polynomials.  Then
\[
U_M(a_1,\dots,a_M)=\sum_{r=0}^M \ii^r e_r(a_1,\dots,a_M).
\]
Consequently, if
\[
S_r:=\sum_{q=0}^r \ii^q e_q(a_1,\dots,a_M),
\qquad 0\le r\le M,
\]
then the successive differences $S_r-S_{r-1}$ are alternately horizontal and vertical segments, and the $r$th segment has signed length $e_r(a_1,\dots,a_M)$.

Moreover,
\[
\widehat U_M(a_1,\dots,a_M)=\exp\Bigl(\ii\sum_{n=1}^M \arctan a_n\Bigr).
\]
\end{proposition}

\begin{proof}
Expanding the product gives
\[
\prod_{n=1}^M (1+\ii a_n)=\sum_{I\subseteq [M]} \ii^{|I|}\prod_{j\in I} a_j.
\]
Grouping by the cardinality $|I|=r$ yields the first formula because
\[
e_r(a_1,\dots,a_M)=\sum_{|I|=r}\prod_{j\in I}a_j.
\]
Since multiplication by $\ii^r$ rotates the real axis through successive quarter-turns, the geometric statement is immediate.

For the normalized identity, note that
\[
\nu(a)=\frac{1+\ii a}{\sqrt{1+a^2}}=
\cos(\arctan a)+\ii\sin(\arctan a)=\e^{\ii\arctan a}.
\]
Multiplying over $n$ proves the claim.
\end{proof}

\begin{proposition}[Infinitesimal orthogonal rotation]\label{prop:derivative-nu}
For $a\in\RR$ one has
\[
\nu'(a)=\frac{\ii}{1+a^2}\,\nu(a).
\]
Consequently, if $a_n:I\to\RR$ are differentiable on an interval $I$, then
\[
\frac{d}{dt}\widehat U_M\bigl(a_1(t),\dots,a_M(t)\bigr)
=\ii\Bigl(\sum_{n=1}^M \frac{a_n'(t)}{1+a_n(t)^2}\Bigr)
\widehat U_M\bigl(a_1(t),\dots,a_M(t)\bigr).
\]
\end{proposition}

\begin{proof}
By Proposition~\ref{prop:finite-orthogonal},
\[
\nu(a)=\e^{\ii\arctan a},
\]
so
\[
\nu'(a)=\ii\frac{1}{1+a^2}\e^{\ii\arctan a}=\frac{\ii}{1+a^2}\,\nu(a).
\]
For the product, differentiate termwise and factor out the common product.
\end{proof}

\begin{proposition}[Equal-step limit]\label{prop:equal-step-limit}
For every $\theta\in\RR$,
\[
\lim_{N\to\infty}\nu\!\left(\frac{\theta}{N}\right)^N=\e^{\ii\theta}.
\]
\end{proposition}

\begin{proof}
By Proposition~\ref{prop:finite-orthogonal},
\[
\nu\!\left(\frac{\theta}{N}\right)^N
=\exp\!\Bigl(\ii N\arctan\!\left(\frac{\theta}{N}\right)\Bigr).
\]
Since $N\arctan(\theta/N)\to \theta$, the conclusion follows.
\end{proof}

\begin{theorem}[Exact Taylor atomization]\label{thm:taylor-atomization}
Let
\[
\Theta(t)=\sum_{n=1}^{\infty}\beta_n t^n
\]
be a real-analytic function on $(-R,R)$ with $\Theta(0)=0$ and $\beta_n\in\RR$.  Then there exists $r\in(0,R)$ such that for every $|t|\le r$ and every $n\ge1$ one has
\[
|\beta_n t^n|<\frac{\pi}{4}.
\]
Define the atomic slopes
\[
\alpha_n(t):=\tan\bigl(\beta_n t^n\bigr),\qquad |t|\le r.
\]
Then the following hold.
\begin{enumerate}[label=(\roman*)]
\item The series $\sum_{n\ge1}|\alpha_n(t)|$ converges uniformly on $[-r,r]$.
\item For every $M\ge1$,
\[
\prod_{n=1}^M \nu\bigl(\alpha_n(t)\bigr)=\exp\!\Bigl(\ii\sum_{n=1}^M\beta_n t^n\Bigr).
\]
\item The infinite product converges uniformly on $[-r,r]$, and
\[
\prod_{n=1}^{\infty}\nu\bigl(\alpha_n(t)\bigr)=\e^{\ii\Theta(t)}.
\]
\item The unnormalized finite products admit the orthogonal-walk expansions
\[
\prod_{n=1}^M (1+\ii\alpha_n(t))=
\sum_{q=0}^M \ii^q e_q\bigl(\alpha_1(t),\dots,\alpha_M(t)\bigr).
\]
Hence the normalized partial products are normalized endpoints of finite orthogonal walks with step lengths $e_q\bigl(\alpha_1(t),\dots,\alpha_M(t)\bigr)$.
\end{enumerate}
\end{theorem}

\begin{proof}
Choose $r\in(0,R)$ so small that
\[
\sum_{n=1}^{\infty}|\beta_n|r^n<\frac{\pi}{4}.
\]
Then for every $n$ and every $|t|\le r$,
\[
|\beta_n t^n|\le \sum_{m=1}^{\infty}|\beta_m|r^m<\frac{\pi}{4}.
\]
On $[-\pi/4,\pi/4]$ one has $|\tan x|\le 2|x|$, so
\[
|\alpha_n(t)|=|\tan(\beta_n t^n)|\le 2|\beta_n|r^n.
\]
Since $\sum_n |\beta_n|r^n<\infty$, the series $\sum_n |\alpha_n(t)|$ converges uniformly on $[-r,r]$.  This proves~(i).

For~(ii), observe that $|\beta_n t^n|<\pi/2$, so
\[
\arctan\alpha_n(t)=\arctan\bigl(\tan(\beta_n t^n)\bigr)=\beta_n t^n.
\]
By Proposition~\ref{prop:finite-orthogonal},
\[
\prod_{n=1}^M \nu\bigl(\alpha_n(t)\bigr)
=\exp\!\Bigl(\ii\sum_{n=1}^M \arctan\alpha_n(t)\Bigr)
=\exp\!\Bigl(\ii\sum_{n=1}^M \beta_n t^n\Bigr).
\]
This proves~(ii).

Since the power series for $\Theta$ converges uniformly on $[-r,r]$, the partial sums
\[
\Theta_M(t):=\sum_{n=1}^M\beta_n t^n
\]
converge uniformly to $\Theta(t)$.  By~(ii),
\[
\prod_{n=1}^M \nu\bigl(\alpha_n(t)\bigr)=\e^{\ii\Theta_M(t)},
\]
so the partial products converge uniformly to $\e^{\ii\Theta(t)}$.  This proves~(iii).

Part~(iv) is Proposition~\ref{prop:finite-orthogonal} applied to the finite sequence $\alpha_1(t),\dots,\alpha_M(t)$.
\end{proof}

\begin{corollary}[Two exact walk models for the same phase]\label{cor:two-walks}
Under the hypotheses of Theorem~\ref{thm:taylor-atomization}, the same phase $\e^{\ii\Theta(t)}$ admits two exact representations.
\begin{enumerate}[label=(\alph*)]
\item The \emph{Taylor walk}
\[
\e^{\ii\Theta(t)}=\sum_{q=0}^{\infty}\frac{\bigl(\ii\Theta(t)\bigr)^q}{q!},
\]
whose step directions cycle with period $4$ and whose step lengths are $|\Theta(t)|^q/q!$.
\item The \emph{atomic Tobley walk}, obtained as the uniform limit of the normalized endpoints of the finite orthogonal walks attached to the atomic slopes $\alpha_n(t)=\tan(\beta_n t^n)$.
\end{enumerate}
\end{corollary}

\begin{proof}
Part~(a) is the Taylor expansion of the exponential function.  Part~(b) is Theorem~\ref{thm:taylor-atomization}(iii)--(iv).
\end{proof}

\section{Taylor atomization of nonvanishing holomorphic germs}

\begin{theorem}[Local phase expansion]\label{thm:nonvanishing-germ}
Let $g$ be holomorphic on a neighborhood of $0\in\CC$, and assume $g(0)\neq0$.  Then there exist $\varepsilon>0$ and a unique holomorphic function
\[
h(z)=\sum_{n=1}^{\infty} c_n z^n
\]
on $|z|<\varepsilon$ such that
\[
g(z)=g(0)\e^{h(z)},
\qquad h(0)=0.
\]
Fix a unit direction $\nu\in\Sph$, and define
\[
\Theta_{g,\nu}(t):=\Ima\,h(t\nu)
=\sum_{n=1}^{\infty}\beta_{g,\nu,n} t^n,
\qquad
\beta_{g,\nu,n}:=\Ima\bigl(c_n\nu^n\bigr).
\]
Then, for $0<t<\varepsilon$,
\[
\ph\bigl(g(t\nu)\bigr)=\ph\bigl(g(0)\bigr)\e^{\ii\Theta_{g,\nu}(t)}.
\]
Moreover, after shrinking $\varepsilon$ if necessary, Theorem~\ref{thm:taylor-atomization} applies to $\Theta_{g,\nu}$, so the local phase is represented by a canonical atomic sequence
\[
\alpha_{g,\nu,n}(t):=\tan\bigl(\beta_{g,\nu,n}t^n\bigr).
\]
\end{theorem}

\begin{proof}
Because $g(0)\neq0$, after shrinking the neighborhood we may assume that $g$ never vanishes on the disc $|z|<\varepsilon$.  Since the disc is simply connected, the nonvanishing holomorphic function $g(z)/g(0)$ admits a unique holomorphic logarithm $h$ with $h(0)=0$.  Expanding $h$ into its Taylor series gives
\[
g(z)=g(0)\e^{h(z)},
\qquad
h(z)=\sum_{n=1}^{\infty} c_n z^n.
\]
Evaluating at $z=t\nu$ yields
\[
g(t\nu)=g(0)\exp\!\Bigl(\sum_{n=1}^{\infty} c_n\nu^n t^n\Bigr).
\]
Taking moduli gives
\[
|g(t\nu)|=|g(0)|\exp\!\Bigl(\sum_{n=1}^{\infty}\Rea(c_n\nu^n)t^n\Bigr),
\]
while dividing by the modulus yields
\[
\ph\bigl(g(t\nu)\bigr)=\ph\bigl(g(0)\bigr)\exp\!\Bigl(\ii\sum_{n=1}^{\infty}\Ima(c_n\nu^n)t^n\Bigr)
=\ph\bigl(g(0)\bigr)\e^{\ii\Theta_{g,\nu}(t)}.
\]
The last claim is exactly Theorem~\ref{thm:taylor-atomization} applied to the real-analytic function $\Theta_{g,\nu}$.
\end{proof}

\begin{proposition}[The first two Taylor atoms]\label{prop:first-two-atoms}
In the situation of Theorem~\ref{thm:nonvanishing-germ}, write
\[
\frac{g(z)}{g(0)}=1+a_1 z+a_2 z^2+O(z^3)
\qquad (z\to0).
\]
Then
\[
c_1=a_1,
\qquad
c_2=a_2-\frac12 a_1^2,
\]
and therefore
\[
\beta_{g,\nu,1}=\Ima(\nu a_1),
\qquad
\beta_{g,\nu,2}=\Ima\!\left(\nu^2\Bigl(a_2-\frac12 a_1^2\Bigr)\right).
\]
\end{proposition}

\begin{proof}
Since $\log(1+w)=w-\frac12 w^2+O(w^3)$ as $w\to0$, applying this with
\[
w=a_1 z+a_2 z^2+O(z^3)
\]
gives
\[
h(z)=a_1 z+\Bigl(a_2-\frac12 a_1^2\Bigr)z^2+O(z^3).
\]
Comparing coefficients with $h(z)=\sum_{n\ge1} c_n z^n$ proves the formula for $c_1$ and $c_2$.  The formulas for the $\beta$'s are immediate from the definition $\beta_{g,\nu,n}=\Ima(c_n\nu^n)$.
\end{proof}

\section{Local phase packets of the non-trivial zeta zeros}

\begin{definition}
Let $\rho\in Z_\xi$ be a non-trivial zero of multiplicity
\[
m:=m_\rho=\ord_\rho(\xi).
\]
Define the leading coefficient
\[
\lambda_\rho:=\frac{\xi^{(m)}(\rho)}{m!}\in\CC^\times.
\]
Then the reduced local germ at $\rho$ is
\[
\psi_\rho(z):=\frac{\xi(\rho+z)}{\lambda_\rho z^m}.
\]
By construction $\psi_\rho(0)=1$.
\end{definition}

\begin{theorem}[Local phase packet at a zeta zero]\label{thm:local-phase-packet}
Let $\rho\in Z_\xi$ have multiplicity $m$, and let $\nu\in\Sph$ be a unit direction.  Then there exists $\varepsilon>0$ and a unique holomorphic function
\[
h_\rho(z)=\sum_{n=1}^{\infty} c_{\rho,n} z^n
\qquad (|z|<\varepsilon)
\]
such that
\[
\xi(\rho+z)=\lambda_\rho z^m \e^{h_\rho(z)},
\qquad h_\rho(0)=0.
\]
Define the local phase series
\[
\Theta_{\rho,\nu}(t):=\Ima\,h_\rho(t\nu)
=\sum_{n=1}^{\infty}\beta_{\rho,\nu,n} t^n,
\qquad
\beta_{\rho,\nu,n}:=\Ima\bigl(c_{\rho,n}\nu^n\bigr).
\]
Then, for $0<t<\varepsilon$,
\begin{align*}
|\xi(\rho+t\nu)|
&=t^m |\lambda_\rho| \exp\!\Bigl(\sum_{n=1}^{\infty}\Rea(c_{\rho,n}\nu^n)t^n\Bigr),\\[4pt]
\ph\bigl(\xi(\rho+t\nu)\bigr)
&=\nu^m\ph(\lambda_\rho)\e^{\ii\Theta_{\rho,\nu}(t)}.
\end{align*}
After shrinking $\varepsilon$ if necessary, the phase factor admits the exact atomic representation
\[
\e^{\ii\Theta_{\rho,\nu}(t)}
=\prod_{n=1}^{\infty}\nu\bigl(\alpha_{\rho,\nu,n}(t)\bigr),
\qquad
\alpha_{\rho,\nu,n}(t):=\tan\bigl(\beta_{\rho,\nu,n}t^n\bigr),
\]
and its finite truncations are normalized endpoints of the orthogonal walks
\[
W_{\rho,\nu,M}(t):=
\sum_{q=0}^M \ii^q e_q\bigl(\alpha_{\rho,\nu,1}(t),\dots,\alpha_{\rho,\nu,M}(t)\bigr).
\]
\end{theorem}

\begin{proof}
Because $\xi(\rho+z)$ has a zero of exact order $m$ at $z=0$, one has a factorization
\[
\xi(\rho+z)=\lambda_\rho z^m \psi_\rho(z)
\]
with $\psi_\rho$ holomorphic and $\psi_\rho(0)=1$.  Shrinking $\varepsilon$ if necessary, we may assume that $\psi_\rho$ never vanishes on $|z|<\varepsilon$.  Since the disc is simply connected, there exists a unique holomorphic logarithm $h_\rho$ with $h_\rho(0)=0$ and
\[
\psi_\rho(z)=\e^{h_\rho(z)}.
\]
This proves the factorization
\[
\xi(\rho+z)=\lambda_\rho z^m \e^{h_\rho(z)}.
\]
Now evaluate at $z=t\nu$ with $0<t<\varepsilon$.  Because $|\nu|=1$,
\[
\xi(\rho+t\nu)=\lambda_\rho t^m \nu^m \exp\!\Bigl(\sum_{n=1}^{\infty} c_{\rho,n}\nu^n t^n\Bigr).
\]
Taking moduli yields the first displayed formula, and dividing by the modulus yields
\[
\ph\bigl(\xi(\rho+t\nu)\bigr)=\nu^m\ph(\lambda_\rho)\exp\!\Bigl(\ii\sum_{n=1}^{\infty}\Ima(c_{\rho,n}\nu^n)t^n\Bigr)
=\nu^m\ph(\lambda_\rho)\e^{\ii\Theta_{\rho,\nu}(t)}.
\]
The atomic representation and the orthogonal-walk formula are exactly Theorem~\ref{thm:taylor-atomization} applied to the real-analytic function $\Theta_{\rho,\nu}$.
\end{proof}

\begin{corollary}[The first local phase atoms of a zeta zero]\label{cor:zeta-first-atoms}
Let $\rho\in Z_\xi$ have multiplicity $m$.  Then
\[
\psi_\rho(z)=1+a_{\rho,1} z+a_{\rho,2} z^2+O(z^3),
\]
where
\[
a_{\rho,1}=\frac{\xi^{(m+1)}(\rho)}{(m+1)\xi^{(m)}(\rho)},
\qquad
a_{\rho,2}=\frac{\xi^{(m+2)}(\rho)}{(m+2)(m+1)\xi^{(m)}(\rho)}.
\]
Hence, for every unit direction $\nu$,
\[
\beta_{\rho,\nu,1}=\Ima\!\left(\nu\frac{\xi^{(m+1)}(\rho)}{(m+1)\xi^{(m)}(\rho)}\right),
\]
and
\[
\beta_{\rho,\nu,2}=\Ima\!\left(\nu^2\Biggl(\frac{\xi^{(m+2)}(\rho)}{(m+2)(m+1)\xi^{(m)}(\rho)}-
\frac12\Bigl(\frac{\xi^{(m+1)}(\rho)}{(m+1)\xi^{(m)}(\rho)}\Bigr)^2\Biggr)\right).
\]
In particular, if $\rho$ is simple and lies on the critical line, then the first normal and tangent atoms are
\[
\beta_{\rho,1,1}=\Ima\!\left(\frac{\xi''(\rho)}{2\xi'(\rho)}\right),
\qquad
\beta_{\rho,\ii,1}=\Rea\!\left(\frac{\xi''(\rho)}{2\xi'(\rho)}\right).
\]
\end{corollary}

\begin{proof}
Since
\[
\xi(\rho+z)=\sum_{k=m}^{\infty}\frac{\xi^{(k)}(\rho)}{k!}z^k,
\]
dividing by $\lambda_\rho z^m$ gives
\[
\psi_\rho(z)=1+
\frac{\xi^{(m+1)}(\rho)}{(m+1)\xi^{(m)}(\rho)}z+
\frac{\xi^{(m+2)}(\rho)}{(m+2)(m+1)\xi^{(m)}(\rho)}z^2+O(z^3).
\]
Now apply Proposition~\ref{prop:first-two-atoms}.  If $m=1$ and $\nu=\ii$, then
\[
\beta_{\rho,\ii,1}=\Ima\!\left(\ii\frac{\xi''(\rho)}{2\xi'(\rho)}\right)=\Rea\!\left(\frac{\xi''(\rho)}{2\xi'(\rho)}\right).
\]
\end{proof}

\begin{proposition}[Packet symmetry of the local phase data]\label{prop:packet-symmetry-local}
Let $\rho\in Z_\xi$ have multiplicity $m$.  Then
\[
\lambda_{1-\rho}=(-1)^m\lambda_\rho,
\qquad
\lambda_{\bar\rho}=\overline{\lambda_\rho}.
\]
Moreover, after shrinking the common radius of definition if necessary, the reduced local germs satisfy
\[
\psi_{1-\rho}(z)=\psi_\rho(-z),
\qquad
\psi_{\bar\rho}(z)=\overline{\psi_\rho(\bar z)}.
\]
Consequently, the logarithmic phase series satisfy
\[
\Theta_{1-\rho,-\nu}(t)=\Theta_{\rho,\nu}(t),
\qquad
\Theta_{\bar\rho,\bar\nu}(t)=-\Theta_{\rho,\nu}(t),
\]
and therefore the atomic slopes transform by
\[
\alpha_{1-\rho,-\nu,n}(t)=\alpha_{\rho,\nu,n}(t),
\qquad
\alpha_{\bar\rho,\bar\nu,n}(t)=-\alpha_{\rho,\nu,n}(t).
\]
In particular,
\[
|\lambda_{\rho}|=|\lambda_{\bar\rho}|=|\lambda_{1-\rho}|=|\lambda_{1-\bar\rho}|,
\]
so the local jet weight
\[
\omega(P_\rho):=|\lambda_\rho|
\]
is a well-defined positive invariant of the packet $P_\rho$.
\end{proposition}

\begin{proof}
Because $\xi(1-s)=\xi(s)$, differentiating $m$ times gives
\[
(-1)^m\xi^{(m)}(1-s)=\xi^{(m)}(s).
\]
Setting $s=\rho$ yields
\[
\xi^{(m)}(1-\rho)=(-1)^m\xi^{(m)}(\rho),
\]
hence $\lambda_{1-\rho}=(-1)^m\lambda_\rho$.  Likewise, $\overline{\xi(\bar s)}=\xi(s)$ implies
\[
\xi^{(m)}(\bar\rho)=\overline{\xi^{(m)}(\rho)},
\]
whence $\lambda_{\bar\rho}=\overline{\lambda_\rho}$.

Next,
\[
\xi(1-\rho+z)=\xi(\rho-z)=\lambda_\rho(-z)^m\psi_\rho(-z)=\lambda_{1-\rho} z^m\psi_\rho(-z),
\]
so by the definition of $\psi_{1-\rho}$ one has $\psi_{1-\rho}(z)=\psi_\rho(-z)$.  Similarly,
\[
\xi(\bar\rho+z)=\overline{\xi(\rho+\bar z)}
=\overline{\lambda_\rho \bar z^m \psi_\rho(\bar z)}
=\lambda_{\bar\rho} z^m\overline{\psi_\rho(\bar z)},
\]
so $\psi_{\bar\rho}(z)=\overline{\psi_\rho(\bar z)}$.

Let $h_\rho$ be the logarithm used in Theorem~\ref{thm:local-phase-packet}.  Since $\psi_{1-\rho}(0)=1$ and $\psi_\rho(0)=1$, uniqueness of the logarithm normalized by value $0$ at the origin gives
\[
h_{1-\rho}(z)=h_\rho(-z),
\qquad
h_{\bar\rho}(z)=\overline{h_\rho(\bar z)}.
\]
Evaluating at $z=-t\nu$ and $z=t\bar\nu$ shows
\[
\Theta_{1-\rho,-\nu}(t)=\Ima\,h_{1-\rho}(-t\nu)=\Ima\,h_\rho(t\nu)=\Theta_{\rho,\nu}(t),
\]
and
\[
\Theta_{\bar\rho,\bar\nu}(t)=\Ima\,\overline{h_\rho(t\nu)}=-\Ima\,h_\rho(t\nu)=-\Theta_{\rho,\nu}(t).
\]
The slope formulas follow because $\tan$ is odd.  The weight invariance is immediate from the transformation law for the $\lambda$'s.
\end{proof}

\begin{theorem}[Geometric image of a non-trivial zero]\label{thm:final}
Let $\rho$ be a non-trivial zero of $\zeta$, let $m=m_\rho$, and let
\[
z_\rho:=\Phi(\rho)=\ii\Bigl(\rho-\frac12\Bigr).
\]
Then the geometric image of $\rho$ in the split-zero formalism is the following packet object.
\begin{enumerate}[label=(\roman*)]
\item The reduced coefficient is the fixed Boolean pair
\[
A=\Fix_S(-1)=\{0,\taup\}\cong\BB,
\]
and the multiplicity coefficient object is the Boolean cube $A^m$.
\item The packet base is the centered spectral packet
\[
\Phi(P_\rho)=\{z_\rho,-z_\rho,\bar z_\rho,-\bar z_\rho\}.
\]
\item The packet carries the well-defined positive weight
\[
\omega(P_\rho)=\left|\frac{\xi^{(m)}(\rho)}{m!}\right|.
\]
\item For every unit direction $\nu\in\Sph$ there is a canonical real-analytic phase series
\[
\Theta_{\rho,\nu}(t)=\sum_{n=1}^{\infty}\beta_{\rho,\nu,n}t^n
\]
and a canonical atomic slope sequence
\[
\alpha_{\rho,\nu,n}(t)=\tan\bigl(\beta_{\rho,\nu,n}t^n\bigr)
\qquad (0<t\ll1),
\]
such that
\[
\ph\bigl(\xi(\rho+t\nu)\bigr)=\nu^m\ph(\lambda_\rho)
\prod_{n=1}^{\infty}\frac{1+\ii\alpha_{\rho,\nu,n}(t)}{\sqrt{1+\alpha_{\rho,\nu,n}(t)^2}}.
\]
The $M$th truncation is the normalized endpoint of the orthogonal walk
\[
W_{\rho,\nu,M}(t)=\sum_{q=0}^M \ii^q e_q\bigl(\alpha_{\rho,\nu,1}(t),\dots,\alpha_{\rho,\nu,M}(t)\bigr).
\]
\item Under the packet symmetries $\rho\mapsto1-\rho$ and $\rho\mapsto\bar\rho$, the phase series and atomic slope sequences transform exactly as in Proposition~\ref{prop:packet-symmetry-local}.  Thus the packet carries a well-defined local phase geometry.
\item If $\rho=\tfrac12+\ii\gamma$ lies on the critical line, then
\[
z_\rho=-\gamma\in\RR
\]
and the packet determines the primitive $\zeta$-cycle length
\[
\Lambda_\rho=\frac{2\pi}{|\gamma|}=\frac{2\pi}{|z_\rho|}.
\]
\end{enumerate}
Equivalently, a non-trivial zeta zero is represented geometrically by a weighted Boolean spectral packet together with its canonical Taylor-atomic phase sequence.
\end{theorem}

\begin{proof}
Item~(i) is Proposition~\ref{prop:boolean-pair} together with Proposition~\ref{prop:F1-zero}.  Item~(ii) is Proposition~\ref{prop:packet-coordinate}.  Item~(iii) is the weight invariance proved in Proposition~\ref{prop:packet-symmetry-local}.  Item~(iv) is Theorem~\ref{thm:local-phase-packet}.  Item~(v) is Proposition~\ref{prop:packet-symmetry-local}.  Item~(vi) is Corollary~\ref{cor:critical-length}.
\end{proof}

\begin{remark}
The theorem separates three levels of structure.
\begin{enumerate}[label=(\arabic*)]
\item The pointwise split-zero coefficient $A$ is constant and therefore too coarse.
\item The Boolean cube $A^m$ and the packet coordinate $\Phi(P_\rho)$ record multiplicity and global symmetry.
\item The Taylor-atomic phase sequence records the genuine local analytic geometry of the zero.
\end{enumerate}
Thus the local orthogonal-walk packet is the first level at which the non-trivial zeros become visible inside the split-zero formalism.
\end{remark}

\begin{thebibliography}{99}

\bibitem{ConnesConsaniCycles}
A.~Connes and C.~Consani,
\emph{Hochschild homology, trace map and $\zeta$-cycles},
in \emph{Cyclic Cohomology at 40: Achievements and Future Prospects},
Proc. Sympos. Pure Math. \textbf{105}, American Mathematical Society, Providence, RI, 2023, pp.~83--110.

\bibitem{Tobley}
J.~Tobley,
\emph{Unpublished correspondence on orthogonal-walk phase decompositions, Taylor-series atomization, and affine cyclotomic charts},
May 2026.  Communicated in the project source transcripts supplied with the manuscript materials.

\end{thebibliography}

\end{document}
